% Generated draft: requires mathematical and source review.
\documentclass[11pt]{article}
\usepackage[margin=25mm]{geometry}
\usepackage{fontspec}
\setmainfont{Latin Modern Roman}
\usepackage{amsmath,amssymb,mathtools,unicode-math}
\setmathfont{Latin Modern Math}
\usepackage{enumitem,needspace,xurl,hyperref}
\setlength{\parindent}{0pt}
\setlength{\parskip}{5pt}
\setlength{\emergencystretch}{3em}
\newcommand{\N}{\mathbb{N}}
\newcommand{\Z}{\mathbb{Z}}
\newcommand{\Q}{\mathbb{Q}}
\newcommand{\R}{\mathbb{R}}
\newcommand{\C}{\mathbb{C}}
\newcommand{\F}{\mathbb{F}}
\newcommand{\A}{\mathbb{A}}
\newcommand{\PP}{\mathbb{P}}
\newcommand{\E}{\mathbb{E}}
\newcommand{\eps}{\varepsilon}
\newcommand{\dd}{\,\mathrm{d}}
\begin{document}
\section*{UnsolvedMath ID: 4800025}
\textbf{Source problem code:} catalogue entry
\subsection{Definitions and mathematical statement}
A Hardy field is a field of germs at \(+\infty\) of real-valued eventually differentiable functions that is closed under differentiation. Let \(\ell\ge1\) and let \(a_1,\ldots,a_\ell\) belong to one Hardy field, with polynomial growth: for each \(i\), \(|a_i(t)|\le C_i t^{k_i}\) for all sufficiently large \(t\), for some \(C_i>0\) and \(k_i\ge0\). On the nonzero-function real-span interpretation, set
\[
\operatorname{span}^*(\mathcal F)=
\left\{\sum_{i=1}^\ell\lambda_i a_i:\lambda_i\in\R\right\}\setminus\{0\},
\qquad\mathcal F=\{a_1,\ldots,a_\ell\}.
\]
Assume the exact separation condition
\[
\forall a\in\operatorname{span}^*(\mathcal F)\ \forall c\in\R\ \forall p\in\Z[t],
\qquad \lim_{t\to\infty}|a(t)-cp(t)|=+\infty.
\]
This concerns every indicated linear combination and every real multiple of an integer polynomial, not merely each individual function or a failure of polynomial equality.

Writing \([s]=\lfloor s\rfloor\), the usual existence-of-recurrence target is: for every invertible probability-preserving system \((X,\mathcal B,\mu,T)\) and every \(A\in\mathcal B\) with \(\mu(A)>0\), there is a positive integer \(n\), in the common domain of the representatives, such that
\[
\mu\left(A\cap T^{-\lfloor a_1(n)\rfloor}A\cap\cdots
\cap T^{-\lfloor a_\ell(n)\rfloor}A\right)>0.
\]
All powers belong to the same transformation. Invertibility accommodates negative integer parts. The source calls this being good for \(\ell\)-recurrence, but does not supply its full convention; whether arbitrarily large \(n\), positive averaged recurrence, or further nontriviality conditions are included must be checked. Those stronger conclusions are not silently added. The coefficient domain and exclusion of zero in \(\operatorname{span}^*\), and the integer-part convention, likewise require the source definitions.

The supplied research summary attributes an affirmative result to Bergelson--Moreira--Richter (2020). Its abbreviated phrase non-polynomial-growth hypothesis does not replace the displayed separation condition, since the functions here explicitly have polynomial growth.
\subsection{Short English statement}
Integer parts of Hardy-field functions should give simultaneous recurrence for one transformation when every nonzero linear combination diverges away from every real multiple of an integer polynomial. The supplied registry reports a literature solution, with exact recurrence conventions still to check.
\subsection{Sources}
\begin{itemize}
\item Local UnsolvedMath registry, permanent ID \(4800025\), exact source code \texttt{catalogue entry}, ``Multiple ergodic averages --- Problem 25.''
\item Nikos Frantzikinakis, \emph{Some open problems on multiple ergodic averages} (2016), Problem 25; \url{https://arxiv.org/abs/1103.3808}. Source-TeX extraction; recorded access 2026-07-29.
\item Supplied research summary attributes the resolution to Bergelson--Moreira--Richter, \emph{Advances} \textbf{368} (2020), bibliographic locator supplied as ``107--146.'' No article title or further theorem details are supplied; this locator's formatting needs verification.
\end{itemize}
\end{document}
